\chapter{Conclusion}
\label{chap:conclusion}
\label{sec:conclusion-summary}

This thesis reconstructs Simplified Jordan Sorting as one executable core with
ordinary-list and heterogeneous finger-tree backends. Both recover their own
sorted output; external certification supplies validity, not the answer.
The resulting artifact connects algorithm reconstruction, a specialized
representation, a source-level argument and a reproducible comparison.

\section{Answers to the Research Questions}
\label{sec:conclusion-questions}
\label{sec:conclusion-rq1}

\textbf{RQ1: reconstruction and correctness.}
The shared Step~1/2/3 loop maintains geometric ownership, exact sibling groups
and sorted prefixes. The argument in \cref{app:backend-argument} explains the
two odd-stage $z_1$ corrections, temporary transfer, reversed ribs and recursive
repair transients. Stable pair identities and child groups survive splitting.
The two frozen input sets passed v2 qualification; all 66,560 formal sample
calls and 1,520 warm-ups returned correct outputs. Finite execution evidence
and the manual proof have different roles; neither is machine verification.

\label{sec:conclusion-rq2}
\textbf{RQ2: observed runtime.}
For the measured configuration retaining access hooks, every input has a
higher tree median. Common-set tree/list ratios have median 2.898 and range
1.114--3.535; the targeted set has median 3.180 and range 3.147--3.213.
These summarize ratios of twenty-batch input medians, not confidence intervals
or every individual call. No crossover is established. The comparison
includes different wrapping, checking and allocation costs. A separate
fifteen-input local diagnostic reduces tree medians by 37.7--52.1\% when access
hooks are removed. Four of fifteen cases, the large heavy constructions,
favor the hook-free tree; the remaining eleven take 1.09--1.68 times the
list time. This does not replace the frozen experiment or establish general speedups.

\label{sec:conclusion-rq3}
\textbf{RQ3: structure and cost.}
One potential pays propagation across mixed updates, giving amortized constant
endpoint insertion and smaller-side logarithmic split. A list-size potential
then pays the split terms over the full run, including bounded ownership and
shared-core work. The linear total uses unit-cost comparisons and machine-word
operations, expected dictionary lookup and amortized growth. Ordinary-list
insertion copying alone has a quadratic lower bound.

The common workload has no two-sided splits; the targeted group contains one
per input with source size 61 and smaller sides 3/7/19. Neither identifies an
isolated split-cost function or broad mixed-workload performance. The practical
slowdown of the measured Python configuration is compatible with the structural
bound. Historical rollback results remain separate and non-causal.

\section{Contributions and Further Work}
\label{sec:conclusion-contributions}
\label{sec:conclusion-future-work}
\label{sec:conclusion-closing}

The contribution is an inspectable implementation and argument with explicit
source, input, timing and statistical boundaries. Independent review of the
finite-parent locality extension and mixed-operation proof would strengthen
confidence in the argument. Future
work should qualify a production hook-free configuration and broaden split
histories and environments under predeclared rules. The post-review diagnostic
and 40-input seeded correctness supplement are bounded first steps.
Such work should preserve the present data regardless of
its speed outcome. Recognition, invalid-input behavior, error correction and
geometric applications require separate specifications. The thesis therefore
reports a structural bound, finite correctness evidence and measured runtime
as distinct conclusions.
